\begin{frontmatter}

\title{Frequent Losers as Cover Bidders in Public Procurement\texorpdfstring{\footnote{This version: March 2026.}}{}}

\author[insper]{Darcio Genicolo-Martins\corref{cor1}}
\ead{darciogm1@insper.edu.br}
\author[insper]{Paulo Furquim de Azevedo}
\affiliation[insper]{organization={INSPER},
  city={S\~ao Paulo}, country={Brazil}}
\cortext[cor1]{Corresponding author.}

\begin{abstract} % REVISION-V6: abstract rewrite — RAND reframe
Cover bidding---the deployment of sham bidders to simulate competition
in rigged procurement auctions---is a widespread but difficult-to-detect
cartel practice. We develop a structural model of cover-bidder deployment
that predicts the optimal number, bid distribution, and market selection
of cover bidders as a function of genuine competition, reference prices,
and detection probability. Estimating the model on 40 million bids from
Brazilian public procurement (2009--2019), we identify 2,735
cover-bidder candidates (``frequent losers'': firms that never win yet
participate abnormally often). The structural estimates imply a cartel
markup consistent with OLS (6.4\%) and within the range of cartel
overcharge estimates in the literature. Validation against competition
authority convictions shows the model-implied detection probability
significantly predicts actual enforcement outcomes. We derive the
welfare-maximizing detection threshold as a function of enforcement
capacity, providing a practical screening tool for competition
authorities.
\end{abstract}

\end{frontmatter}

\noindent\textbf{Keywords:} bid rigging, cover bidding, public procurement, cartel screening, frequent losers,\\ instrumental variables

\noindent\textbf{JEL No.:} D44, D73, H57, K21, L41

\bigskip
